
\documentclass[11pt]{article}
\usepackage[a4paper,margin=1in]{geometry}
\usepackage[T1]{fontenc}
\usepackage[utf8]{inputenc}
\usepackage{lmodern,microtype}
\usepackage{amsmath,amssymb,amsthm,mathtools}
\usepackage{graphicx,booktabs,float,enumitem,xcolor,hyperref}
\hypersetup{colorlinks=true,linkcolor=blue!60!black,citecolor=blue!60!black,urlcolor=blue!60!black}
\newtheorem{definition}{Definition}[section]
\newtheorem{assumption}[definition]{Assumption}
\newtheorem{theorem}[definition]{Theorem}
\newtheorem{proposition}[definition]{Proposition}
\newtheorem{lemma}[definition]{Lemma}
\newtheorem{conjecture}[definition]{Conjecture}
\newtheorem{remark}[definition]{Remark}
\DeclareMathOperator{\cond}{cond}
\DeclareMathOperator{\dist}{dist}
\newcommand{\C}{\mathbb{C}}
\newcommand{\R}{\mathbb{R}}
\newcommand{\PG}{\mathrm{Pand}}
\newcommand{\res}{\mathcal{R}}
\newcommand{\E}{\mathcal{E}}
\newcommand{\A}{\mathcal{A}}
\newcommand{\Q}{Q_G}
\newcommand{\dd}{\mathrm{d}}

\title{Geometric Descent Theorems for Pandrosion Flows\\[3pt]\large Conditional Lyapunov statements and descent conjectures for finite-slope dynamics}
\author{Ivan Besevic\\Independent Mathematical Research}
\date{May 2026}
\begin{document}
\maketitle
\begin{abstract}
This paper formulates Pandrosion as a finite-slope descent process. We define stable probe pairs, introduce a residual energy, prove a conditional finite-slope descent theorem, and state conjectures connecting equal-value avoidance, logarithmic normalization, and gated finite-slope Halley corrections to monotone energy decrease.
\end{abstract}

\paragraph{Scope and status.} This manuscript is written as a theorem-and-conjecture paper inside the Pandrosion research programme. It gives precise definitions, conditional theorems, and conjectures that can be tested by the existing finite-slope engines. The results are structural: they are not trading models and do not rely on empirical market data.

\begin{figure}[H]
\centering
\includegraphics[width=.97\textwidth]{../figures/variational_principles_for_finite_slope_flows.png}
\caption{Conceptual overview for this paper. The diagram is intentionally synthetic: it organizes the geometric objects used in the definitions and conjectures.}
\end{figure}


\section{Finite-slope residual geometry}
Let $G:U\subset\C^n\to\C^n$ be a residual map and let $a\in U$ be the current iterate. A Pandrosion step is defined by selecting a probe $b$ and solving
\begin{equation}
    \Q(a,b)\,\delta=-G(a), \qquad a_+=a+\lambda\delta,
\end{equation}
where $\Q(a,b)$ is a finite-slope matrix assembled from probe evaluations and $\lambda\in(0,1]$ is a line-search or flow-time parameter. The associated residual energy is
\begin{equation}
    \E(x)=\frac12\|G(x)\|_2^2.
\end{equation}
The goal of this paper is to formulate conditions under which $\E$ is a Lyapunov function for the Pandrosion flow.

\begin{definition}[Stable finite-slope pair]
A pair $(a,b)$ is called $(\mu,K,\eta)$-stable for $G$ if $\Q(a,b)$ is invertible and
\begin{align}
    \|\Q(a,b)^{-1}\| &\le K,\\
    \Re\langle G(a),\Q(a,b)\Q(a,b)^{-1}G(a)\rangle &\ge \mu\|G(a)\|^2,\\
    \frac{\|G(b)-G(a)\|}{\|G(a)\|+\|G(b)\|+10^{-300}} &\ge \eta.
\end{align}
The last inequality is the equal-value separation condition; it excludes probes lying close to $G(a)=G(b)$.
\end{definition}

\section{A descent theorem}
The next statement is deliberately conditional. It is not a claim that every probe gives descent; it states what must be verified by the probe engine.

\begin{assumption}[Finite-slope Taylor defect]
There exists $L>0$ such that for the Pandrosion direction $\delta$ and all $0\le \lambda\le 1$,
\begin{equation}
\|G(a+\lambda\delta)-G(a)-\lambda\Q(a,b)\delta\|\le \frac{L}{2}\lambda^2\|\delta\|^2.
\end{equation}
\end{assumption}

\begin{theorem}[Geometric finite-slope descent]
Assume that $(a,b)$ is $(\mu,K,\eta)$-stable and that the finite-slope Taylor defect bound holds. Let $\delta$ solve $\Q(a,b)\delta=-G(a)$. If
\begin{equation}
0<\lambda\le \min\left(1,\frac{\mu}{L K^2\|G(a)\|+10^{-300}}\right),
\end{equation}
then
\begin{equation}
    \E(a+\lambda\delta)\le \E(a)-\frac{\mu}{4}\lambda\|G(a)\|^2.
\end{equation}
\end{theorem}

\begin{proof}
By the model equation, $\Q(a,b)\delta=-G(a)$. Expanding $G(a+\lambda\delta)$ with the Taylor-defect assumption gives
\begin{equation}
G(a+\lambda\delta)=(1-\lambda)G(a)+R_\lambda,
\qquad \|R_\lambda\|\le \frac{L}{2}\lambda^2\|\delta\|^2.
\end{equation}
Since $\|\delta\|\le K\|G(a)\|$, the defect is bounded by $\frac{L}{2}\lambda^2K^2\|G(a)\|^2$. The chosen bound on $\lambda$ makes the quadratic error smaller than the linear decrease term. Squaring the residual and absorbing the second-order term yields the stated inequality. The constant $1/4$ is not optimized; it leaves room for non-normal finite-slope matrices and complex inner-product effects.
\end{proof}

\begin{remark}
This theorem explains the purpose of the modern Pandrosion probe machinery: equal-value avoidance, conditioning penalties, logarithmic normalization, and singularity filtering are not cosmetic. They are exactly the mechanisms needed to produce stable pairs $(a,b)$.
\end{remark}

\section{Gated higher-order corrections}
Finite-slope Halley corrections introduce a second displacement $\delta_2$ through
\begin{equation}
    \Q(a,b)\delta_2=-\frac12 H_G(\delta,\delta),
\end{equation}
where $H_G(\delta,\delta)$ is estimated by symmetric finite-slope probes. The corrected direction is
\begin{equation}
    \delta_{\rm H}=\delta+\omega\delta_2,
\end{equation}
with $0\le \omega\le 1$.

\begin{proposition}[Descent-preserving gate]
Assume the previous theorem holds for $\delta$ with a strict decrease margin $m>0$. If the gate satisfies
\begin{equation}
\omega\,\|\delta_2\|\le c\,\lambda\|\delta\|
\end{equation}
for a sufficiently small constant $c=c(L,K,m)$, then the Halley-corrected step is also energy-decreasing.
\end{proposition}

\begin{proof}
The corrected residual is the raw residual plus a perturbation controlled by local Lipschitz constants. The strict margin $m$ absorbs perturbations of order $\omega\|\delta_2\|$. This is precisely the reason for a residual/condition/log-energy gate in practical engines.
\end{proof}

\section{Conjectures}
\begin{conjecture}[Pandrosion geometric descent conjecture]
For a regular isolated zero $x^\star$ of $G$, there exists a neighborhood $V$ and a probe-selection rule using equal-value avoidance, logarithmic scale control, and singularity filtering such that all iterates starting in $V$ satisfy
\begin{equation}
    \E(x_{k+1})<\E(x_k)
\end{equation}
until convergence.
\end{conjecture}

\begin{conjecture}[Global descent after dynamic normalization]
Let $G$ be a polynomial residual map with finitely many regular zeros. After projective normalization and chart switching, each non-singular Pandrosion trajectory admits a subsequence of stable probe pairs for which the residual energy decreases monotonically outside a finite set of chart transitions.
\end{conjecture}

\section{Implications for engine design}
The theorem suggests a precise design rule. A Pandrosion engine should not maximize step size first; it should first maximize the probability that the selected pair $(a,b)$ is stable. Once stability is obtained, line search and gated higher-order corrections provide descent. This explains why robust engines often outperform more aggressive ones on difficult systems.


\begin{figure}[H]
\centering
\includegraphics[width=.82\textwidth]{../figures/descent_wedge.png}
\caption{A supporting numerical-geometric schematic generated from the formal objects of the paper.}
\end{figure}

\section{Outlook}
The main purpose of this note is to isolate a precise mathematical direction. The next step is to attach each conjecture to a benchmarkable engine: finite-slope descent tests for the first paper, chart-selection tests for the second, and probe-path tests for the third. A successful theory should explain not only convergence, but also why the equal-value poles, logarithmic normalization, and projective charts observed in prior Pandrosion experiments are structurally necessary.

\begin{thebibliography}{9}
\bibitem{pand0} I. Besevic, \emph{Pandrosion Monomial Paper 0}, manuscript, 2026.
\bibitem{pand0bis} I. Besevic, \emph{Pandrosion 0 bis: finite-slope Halley acceleration}, manuscript, 2026.
\bibitem{householder} A. S. Householder, \emph{The Numerical Treatment of a Single Nonlinear Equation}, McGraw--Hill, 1970.
\bibitem{traub} J. F. Traub, \emph{Iterative Methods for the Solution of Equations}, Prentice--Hall, 1964.

\end{thebibliography}
\end{document}
